\section{BIBLIOGRAFÍA}

\begin{itemize}[label={}, leftmargin=1.5em, itemindent=-1.5em, itemsep=0.5em]

\item Alagic, G., Apon, D., Cooper, D., Dang, Q., Dang, T., Kelsey, J., Lichtinger, J., Liu, Y.-K., Miller, C., Moody, D., Peralta, R., Perlner, R., Robinson, A., \& Smith-Tone, D. (2022). Status report on the third round of the NIST post-quantum cryptography standardization process (NIST IR 8413). National Institute of Standards and Technology. s.l.

\item Alagic, G., Barker, E., Chen, L., Moody, D., Robinson, A., Silberg, H., \& Waller, N. (2025). NIST SP 800-227: Recommendations for key-encapsulation mechanisms. National Institute of Standards and Technology. s.l.

\item Balboa Santander, R. (2006). Encriptación asimétrica híbrida con distribución cuántica de claves (QKD). Tesis de Licenciatura. Universidad Mayor de San Andrés, Carrera de Informática. Bolivia: La Paz.

\item Barros Da Silva, J. V., Gupta, A., Singh, M., \& Singh, I. (2025). QLink: Quantum-safe bridge architecture for blockchain interoperability. Quantum-Hub ICDS Symposium. Penn State University. USA: Pennsylvania.

\item Bolgar, C. (2026). Microsoft's Majorana 2 chip: Better reliability, on track for 2029. Microsoft Source. s.l. \url{https://blogs.microsoft.com/source/}

\item Chaurasia, A., Sharma, A. K., \& Sharma, S. K. (2025). Post quantum cryptography for blockchain security using CRYSTALS-Dilithium authentication. Journal of Discrete Mathematical Sciences \& Cryptography, 28(8), 3121--3132. s.n. s.l. \url{https://doi.org/10.47974/JDMSC-2588}

\item Ezeh, F. E., Oparah, S. O., Gado, P., Gbaraba, S. V., \& Adeleke, A. S. (2025). Designing a post-quantum blockchain voting protocol with zero-knowledge proofs for tamper-resilient electoral infrastructure. International Journal of Multidisciplinary Futuristic Development, 6(2), 71--79. s.n. s.l. \url{https://doi.org/10.54660/IJMFD.2025.6.2.71-79}

\item Farzaliyev, V., Willemson, J., \& Kaasik, J. K. (2021). Improved lattice-based mix-nets for electronic voting. Cybernetica. Estonia: Tallinn.

\item Finogina, T., \& Herranz, J. (2023). On remote electronic voting with both coercion resistance and cast-as-intended verifiability. Journal of Information Security and Applications, 76, 103554. s.n. s.l.

\item Grassi, L., Khovratovich, D., Rechberger, C., Roy, A., \& Schofnegger, M. (2021). Poseidon: A new hash function for zero-knowledge proof systems. Proceedings of the 30th USENIX Security Symposium. s.n. s.l.

\item Khan, P. S. (2025). Quantum-resistant blockchain-based e-voting system with enhanced voter privacy and real-time transparency. International Journal of Computer Application, 15(4), 175--182. s.n. s.l. \url{https://doi.org/10.5281/zenodo.16728461}

\item Lopez, J., Cadena, V., \& Rahman, M. S. (2025). Evaluating post-quantum cryptographic algorithms on resource-constrained devices. arXiv:2507.08312. s.n. s.l.

\item Moody, D. (2026). NIST PQC: The road ahead. Conferencia presentada en ETSI/IQC Quantum Safe Cryptography Conference. s.n. s.l.

\item Mudaliar, K. (2026). Researchers find errors in Microsoft's Majorana 1 findings, casting doubts on quantum claims. Neowin. s.l. \url{https://www.neowin.net}

\item Nina Chura, L. A. (2014). Algoritmo cuántico vs. clásico aplicado a la compresión de información. Tesis de Licenciatura. Universidad Mayor de San Andrés, Carrera de Informática. Bolivia: La Paz.

\item Poudel, A., Poudel, U., Aryal, D., \& Nepal, A. (2024). A quantum-secure and blockchain-integrated e-voting framework with identity validation. School of CS and Engineering, Vellore Institute of Technology. India: Vellore.

\item Reddy, T. K., \& Tschofenig, H. (2026). Post-quantum cryptography recommendations for TLS-based applications. IETF Internet-Draft draft-ietf-uta-pqc-app-02. s.n. s.l.

\item Roux-Langlois, A. (2021). Post-quantum online voting scheme. Financial Cryptography and Data Security Workshops. s.n. s.l.

\item Silde, T. (2021). Electronic voting using lattice-based commitments and verifiable encryption. Aarhus Crypto Seminar. s.n. Denmark: Aarhus.

\item Tortola, D., Pelosi, A., Russo, G. G., Mori, P., \& Ricci, L. (2024). zkSNARKs libraries for blockchains: A comparative study. 6th Distributed Ledger Technology Workshop (DLT2024). s.n. s.l.

\item Wabartha, C. (2025). Post-quantum secure instantiation of the Ordinos e-voting system. Tesis de Maestría. Universität Stuttgart. Germany: Stuttgart.

\item Yu, P. (2025). Quantum-resistant zero-knowledge proof blockchain electronic voting system. Computer Fraud and Security, 2025(02), 958--981. ISSN: 1873-7056. s.n. s.l.

\end{itemize}

% =========================================================
\subsection*{Webgrafía}
% =========================================================

\begin{itemize}[label={}, leftmargin=1.5em, itemindent=-1.5em, itemsep=0.5em]

\item Encryption Consulting. (2026). EU defines clear roadmap for post quantum cryptography transition by 2035. s.n. s.l. Recuperado el 2 de julio de 2026, de \url{https://www.encryptionconsulting.com}

\item FS-ISAC. (2023). Post-quantum cryptography (PQC) working group risk model. TLP White. s.n. s.l. Recuperado el 2 de julio de 2026, de \url{https://www.fsisac.com}

\item G7 Cyber Expert Group. (2026). Statement on advancing a coordinated roadmap for the transition to post-quantum cryptography in the financial sector. s.n. s.l. Recuperado el 2 de julio de 2026, de \url{https://www.g7germany.de}

\item QRAMM. (2024). NIST post-quantum cryptography standards: Complete guide to FIPS 203, 204, 205. s.n. s.l. Recuperado el 2 de julio de 2026, de \url{https://www.qramm.com}

\end{itemize}

% =========================================================
\subsection*{Legislación y estándares técnicos}
% =========================================================

\begin{itemize}[label={}, leftmargin=1.5em, itemindent=-1.5em, itemsep=0.5em]

\item National Institute of Standards and Technology (NIST). (2024). FIPS 203: Module-lattice-based key-encapsulation mechanism standard. U.S. Department of Commerce. s.l. \url{https://doi.org/10.6028/NIST.FIPS.203}

\item National Institute of Standards and Technology (NIST). (2024). FIPS 204: Module-lattice-based digital signature standard. U.S. Department of Commerce. s.l. \url{https://doi.org/10.6028/NIST.FIPS.204}

\item National Institute of Standards and Technology (NIST). (2024). FIPS 205: Stateless hash-based digital signature standard. U.S. Department of Commerce. s.l. \url{https://doi.org/10.6028/NIST.FIPS.205}

\end{itemize}
